% Standalone problem document, generated from the adjacent README.md.
% Compile with XeLaTeX. Mathematical content is maintained in Markdown.
\documentclass[11pt,a4paper]{article}
\usepackage[fontsize=10.5pt]{scrextend}
\usepackage[margin=22mm,headheight=15pt,headsep=9mm,footskip=11mm]{geometry}
\usepackage{fontspec}
\setmainfont{texgyrepagella}[Extension=.otf,UprightFont=*-regular,BoldFont=*-bold,ItalicFont=*-italic,BoldItalicFont=*-bolditalic]
\setsansfont{texgyreheros}[Extension=.otf,UprightFont=*-regular,BoldFont=*-bold,ItalicFont=*-italic,BoldItalicFont=*-bolditalic]
\setmonofont{texgyrecursor}[Extension=.otf,UprightFont=*-regular,BoldFont=*-bold,ItalicFont=*-italic,BoldItalicFont=*-bolditalic,Scale=MatchLowercase]
\usepackage{amsmath,amssymb}
\usepackage{unicode-math}
\setmathfont{texgyrepagella-math.otf}
\usepackage{xcolor,microtype,enumitem,needspace,fancyhdr}
\usepackage{bookmark}
\definecolor{ink}{HTML}{17344B}
\definecolor{accent}{HTML}{197D80}
\definecolor{muted}{HTML}{596773}
\hypersetup{colorlinks=true,linkcolor=accent,urlcolor=accent,pdftitle={IE-25 - Perturbed
MILU conditioning for Neumann problems in two and three
dimensions},pdfauthor={Open Problems in Numerical Linear Algebra}}
\urlstyle{same}
\setlength{\parindent}{0pt}
\setlength{\parskip}{5pt plus 1pt minus 1pt}
\linespread{1.04}
\setlength{\emergencystretch}{3em}
\widowpenalty=10000
\clubpenalty=10000
\displaywidowpenalty=10000
\allowdisplaybreaks[1]
\setlist{nosep,leftmargin=1.5em}
\providecommand{\tightlist}{\setlength{\itemsep}{0pt}\setlength{\parskip}{0pt}}
\providecommand{\pandocbounded}[1]{#1}
\pagestyle{fancy}
\fancyhf{}
\fancyhead[L]{\sffamily\footnotesize\color{muted}OPEN PROBLEMS IN NUMERICAL LINEAR ALGEBRA}
\fancyhead[R]{\sffamily\bfseries\footnotesize\color{accent}IE-25}
\fancyfoot[L]{\parbox[b]{0.9\textwidth}{\sffamily\fontsize{8}{10}\selectfont\color{muted}Editorial ratings; a bounded search cannot certify that a problem remains unsolved.\\Literature check: 2026-09-10}}
\fancyfoot[R]{\sffamily\footnotesize\color{muted}\thepage}
\renewcommand{\headrulewidth}{0.3pt}
\renewcommand{\headrule}{\hbox to\headwidth{\color{accent}\leaders\hrule height \headrulewidth\hfill}}
\makeatletter
\renewcommand{\section}{\Needspace{5\baselineskip}\@startsection{section}{1}{0pt}{12pt plus 2pt minus 2pt}{4pt}{\sffamily\bfseries\large\color{ink}}}
\renewcommand{\subsection}{\Needspace{4\baselineskip}\@startsection{subsection}{2}{0pt}{9pt plus 1pt minus 1pt}{3pt}{\sffamily\bfseries\normalsize\color{ink}}}
\makeatother
\setcounter{secnumdepth}{0}
\begin{document}
{\sffamily\small\color{muted}Linear systems and elimination\par}
\vspace{3pt}
{\raggedright\hyphenpenalty=10000\exhyphenpenalty=10000\sffamily\bfseries\fontsize{21}{24}\selectfont\color{ink}Perturbed
MILU conditioning for Neumann problems in two and three dimensions\par}
\vspace{7pt}
{\sffamily\small\textbf{Difficulty:} challenging\quad\textbf{Importance:} interesting
to the community\par}
{\sffamily\small\bfseries\color{accent}Status: Open\par}
\vspace{4pt}
{\color{accent}\hrule height 0.8pt}
\vspace{6pt}
\textbf{Rating rationale:} Controlling the perturbed factorization
uniformly for singular Neumann systems in two and three dimensions is
challenging. Community impact comes from a concrete preconditioner for
widely used PDE linear solves.

\subsection{Statement}\label{statement}

Fix \(d\in\{2,3\}\) and a bounded connected smooth domain
\(\Omega\subset\mathbb R^d\). On \(h\mathbb Z^d\), retain points whose
cubic control cells \(C_p=p+[-h/2,h/2]^d\) intersect \(\Omega\), ordered
lexicographically as \(1,\ldots,N_h\). For face neighbors set

\[
w_{pq}=\frac{\operatorname{vol}_{d-1}(C_p\cap C_q\cap\Omega)}{h^{d-1}},
\]

and set other weights to zero. Let \(a_p=\sum_qw_{pq}\), \(A_{pp}=a_p\),
\(A_{pq}=-w_{pq}\) for \(p\ne q\), and let \(L\) be the strictly lower
triangular part of \(A\). This is the Purvis--Burkhalter Neumann
discretization, with length fractions in two dimensions and face-area
fractions in three dimensions.

For \(\varepsilon>0\), define
\(E=\operatorname{diag}(e_1,\ldots,e_{N_h})\) recursively by

\[
e_1=a_1,\qquad
e_p=(1+\varepsilon)a_p-
\sum_{\substack{q<p\\w_{pq}>0}}
\frac{w_{pq}}{e_q}\sum_{s>q}w_{qs}\quad(p>1),
\]

and \(M=(L+E)E^{-1}(E+L^T)\). This expresses the source's PMILU
recurrence by neighbor sums, including its explicitly unperturbed first
pivot. Its two-dimensional display is Definition 5; §4 states the
conjecture in both dimensions.

When the pivots are positive, let \(\kappa_+(M^{-1}A)\) be the ratio of
the largest to smallest positive eigenvalue of \(M^{-1/2}AM^{-1/2}\);
the zero eigenvalue from the Neumann nullspace is excluded.

\textbf{Conjecture.} There exist \(c_\Omega,K_\Omega,h_\Omega>0\),
independent of \(h\), such that choosing \(\varepsilon=c_\Omega h^2\)
gives positive pivots and

\[
\kappa_+(M^{-1}A)\le K_\Omega h^{-1}
\qquad(0<h<h_\Omega).
\]

The existence quantifier spells out the source's choice of ``some
moderate constant''; it does not impose an invented interval for that
constant. The authors suggest the practical choice \(c_\Omega=1\).

\subsection{Numerical significance}\label{numerical-significance}

The question concerns a concrete diagonal perturbation of incomplete
factorization for singular Neumann systems. It is separate from relaxed
ILU: the diagonal recurrence differs, and the proposed PMILU bound
includes three dimensions.

\newpage
\subsection{References and status
check}\label{references-and-status-check}

\begin{itemize}
\tightlist
\item
  B. Lee and C. Min, \emph{Optimal preconditioners on solving the
  Poisson equation with Neumann boundary conditions}, J. Comput. Phys.
  \textbf{433} (2021), 110189,
  \href{https://doi.org/10.1016/j.jcp.2021.110189}{DOI}.
  \href{https://math.ewha.ac.kr/~chohong/publications/article_49_copy.pdf}{Author
  text}: §2.1, equation (3); §4, Definition 5 and conjecture, pp.~5--6;
  §6.2 for the three-dimensional experiments.
\item
  G. Hwang et al., \emph{Localized Estimation of Condition Numbers for
  MILU Preconditioners on a Graph}, SIAM J. Numer. Anal. \textbf{64}
  (2026), 1443--1477, \href{https://doi.org/10.1137/24M1722134}{DOI};
  \href{https://arxiv.org/abs/2501.00245}{available preprint}, §2.1 and
  §5, p.~21.
\end{itemize}

On 2026-09-10, checked the original full author text, the later
preprint's positive-definite assumptions and explicit Neumann
future-work statement, its 2026 publication record, and targeted
PMILU/Neumann/title searches including 2025--2026. No resolution was
found. The 2026 journal full text was unavailable; this is a bounded
check using its primary abstract and author preprint. The source's
descriptive phrase ``MILU of \(A+\varepsilon D\)'' differs at the first
pivot from its displayed Definition 5; this entry follows the display.

\subsection{Independent audit ---
2026-09-10}\label{independent-audit-2026-09-10}

Definition 5 and the §4 conjecture in the full Lee--Min author PDF
support the displayed PMILU recurrence, including the first-pivot
exception. The paper proves the rectangular RILU result, not a PMILU
theorem; its smooth-domain PMILU evidence is numerical. The
positive-definite assumptions and Neumann future-work boundary of the
later preprint were checked, along with targeted subsequent searches. No
exact proved subfamily warranting a partial-resolution label, or general
resolution, was located.
\end{document}
